% GENERATED FILE. Edit canonical lessons, then run npm run generate:books.
% canonical-source-hash: fnv1a64:256fe56fb0a3edc6
% canonical-lessons: RU-C339-ubezhdat, RU-C339-pereubezhdat, RU-C339-ubezhdenie, RU-C339-uverennost, RU-C339-somnevatsya

\chapter{Changing minds}
\label{ch:ru-changing-minds}
\hlchaptermodality{339}

\begin{chapteropening}
\textbf{By the end of this chapter:} \emph{I can try to convince someone, change their mind, speak of my convictions and my confidence, and say that I doubt something.}

Complete the last lesson of chapter 339: I can try to convince someone, change their mind, speak of my convictions and my confidence, and say that I doubt something.
\end{chapteropening}

\section[\texorpdfstring{\ru{убеждать} (ubezhdát)}{убеждать (ubezhdát)}]{\ru{убеждать} (ubezhdát) — to convince, to persuade}

\label{lesson:RU-C339-ubezhdat}



\begin{quote}
Before the new one: say the Russian for proof, then the Russian for convincing.
\end{quote}

\subsection*{You'll want to know: \ru{убеждать}}
\textbf{\ru{убеждать}} (\emph{ubezhdát}) — \textquotedblleft{}to convince, to persuade\textquotedblright{}.

\begin{quote}
\ru{Я} \ru{убеждал} \ru{друга} \ru{поехать} \ru{с} \ru{нами}, \ru{но} \ru{он} \ru{не} \ru{захотел}.
\end{quote}

\begin{quote}
\emph{Ya ubezhdál drúga poyékhat s námi, no on ne zakhotél.}
\end{quote}

\begin{quote}
I tried to persuade my friend to come with us, but he did not want to.
\end{quote}

Imperfective; perfective \emph{\ru{убедить}}, which has no form for \emph{\ru{я}} in the future. \emph{\ru{Убеждать} \ru{кого}-\ru{то}} + infinitive, or \emph{\ru{в}} + prepositional.

\begin{cousinweb}
\textbf{\ru{убеждать}} is the imperfective of \emph{\ru{убедить}}, on the root of \emph{\ru{беда}} \textquotedblleft{}trouble\textquotedblright{}, Proto-Slavic \emph{*běda} (EtymDB). The \emph{\ru{жд}} is the Church Slavonic way of softening \emph{\ru{д}}; native Russian softens it to \emph{\ru{ж}}, as in \emph{\ru{вожу}} beside \emph{\ru{водить}}.
\end{cousinweb}

\begin{grammarlens}[title={What you've built}]
One more word for telling, explaining and describing at B1.
\end{grammarlens}

\subsection*{Your turn}
\begin{itemize}
\raggedright
  \item \emph{Say it:} \emph{ubezhdát}
  \item \emph{Say it:} \emph{ubezhdát}, once more
  \item \emph{Say it:} a sentence of your own with \emph{ubezhdát}
  \item \emph{Recall:} say the Russian for proof, then the Russian for convincing, then say \emph{ubezhdát} again
\end{itemize}

\begin{tcolorbox}[breakable,colback=teal!4,colframe=teal!35!black,title={Before you move on}]
What does \ru{убеждать} mean? (\textquotedblleft{}To convince, to persuade\textquotedblright{}.) Say it once more.
\end{tcolorbox}
\section[\texorpdfstring{\ru{переубеждать} (pereubezhdát)}{переубеждать (pereubezhdát)}]{\ru{переубеждать} (pereubezhdát) — to change someone’s mind}

\label{lesson:RU-C339-pereubezhdat}



\begin{quote}
Before the new one: say the Russian for convincing, then the Russian for to convince.
\end{quote}

\subsection*{You'll want to know: \ru{переубеждать}}
\textbf{\ru{переубеждать}} (\emph{pereubezhdát}) — \textquotedblleft{}to change someone’s mind\textquotedblright{}.

\begin{quote}
\ru{Маму} \ru{невозможно} \ru{переубедить}.
\end{quote}

\begin{quote}
\emph{Mámu nevozmózhno pereubedít.}
\end{quote}

\begin{quote}
It is impossible to change Mum’s mind.
\end{quote}

Imperfective; perfective \emph{\ru{переубедить}}. \emph{\ru{Пере}-} adds the sense of over again: to convince someone the other way.

\begin{cousinweb}
\textbf{\ru{переубеждать}} is built from \emph{\ru{пере}-} \textquotedblleft{}over, again\textquotedblright{} and \emph{\ru{убеждать}} \textquotedblleft{}to convince\textquotedblright{}, whose root is that of \emph{\ru{беда}}, Proto-Slavic \emph{*běda} (EtymDB).
\end{cousinweb}

\begin{grammarlens}[title={What you've built}]
One more word for telling, explaining and describing at B1.
\end{grammarlens}

\subsection*{Your turn}
\begin{itemize}
\raggedright
  \item \emph{Say it:} \emph{pereubezhdát}
  \item \emph{Say it:} \emph{pereubezhdát}, once more
  \item \emph{Say it:} a sentence of your own with \emph{pereubezhdát}
  \item \emph{Recall:} say the Russian for convincing, then the Russian for to convince, then say \emph{pereubezhdát} again
\end{itemize}

\begin{tcolorbox}[breakable,colback=teal!4,colframe=teal!35!black,title={Before you move on}]
What does \ru{переубеждать} mean? (\textquotedblleft{}To change someone’s mind\textquotedblright{}.) Say it once more.
\end{tcolorbox}
\section[\texorpdfstring{\ru{убеждение} (ubezhdéniye)}{убеждение (ubezhdéniye)}]{\ru{убеждение} (ubezhdéniye) — a conviction, a firm belief}

\label{lesson:RU-C339-ubezhdenie}



\begin{quote}
Before the new one: say the Russian for to convince, then the Russian for to change someone’s mind.
\end{quote}

\subsection*{You'll want to know: \ru{убеждение}}
\textbf{\ru{убеждение}} (\emph{ubezhdéniye}) — \textquotedblleft{}a conviction, a firm belief\textquotedblright{}.

\begin{quote}
\ru{Он} \ru{не} \ru{ест} \ru{мясо} \ru{по} \ru{убеждению}.
\end{quote}

\begin{quote}
\emph{On ne yest myáso po ubezhdéniyu.}
\end{quote}

\begin{quote}
He does not eat meat out of conviction.
\end{quote}

Neuter. \emph{\ru{По} \ru{убеждению}} — out of conviction. The plural \emph{\ru{убеждения}} means one’s beliefs.

\begin{cousinweb}
\textbf{\ru{убеждение}} is built on \emph{\ru{убедить}} \textquotedblleft{}to convince\textquotedblright{} with the noun suffix \emph{-\ru{ение}}, keeping the Church Slavonic \emph{\ru{жд}}. The root is that of \emph{\ru{беда}}, Proto-Slavic \emph{*běda} (EtymDB).
\end{cousinweb}

\begin{grammarlens}[title={What you've built}]
One more word for telling, explaining and describing at B1.
\end{grammarlens}

\subsection*{Your turn}
\begin{itemize}
\raggedright
  \item \emph{Say it:} \emph{ubezhdéniye}
  \item \emph{Say it:} \emph{ubezhdéniye}, once more
  \item \emph{Say it:} a sentence of your own with \emph{ubezhdéniye}
  \item \emph{Recall:} say the Russian for to convince, then the Russian for to change someone’s mind, then say \emph{ubezhdéniye} again
\end{itemize}

\begin{tcolorbox}[breakable,colback=teal!4,colframe=teal!35!black,title={Before you move on}]
What does \ru{убеждение} mean? (\textquotedblleft{}A conviction, a firm belief\textquotedblright{}.) Say it once more.
\end{tcolorbox}
\section[\texorpdfstring{\ru{уверенность} (uvérennost)}{уверенность (uvérennost)}]{\ru{уверенность} (uvérennost) — confidence, certainty}

\label{lesson:RU-C339-uverennost}



\begin{quote}
Before the new one: say the Russian for to change someone’s mind, then the Russian for a conviction.
\end{quote}

\subsection*{You'll want to know: \ru{уверенность}}
\textbf{\ru{уверенность}} (\emph{uvérennost}) — \textquotedblleft{}confidence, certainty\textquotedblright{}.

\begin{quote}
\ru{Мне} \ru{не} \ru{хватает} \ru{уверенности} \ru{в} \ru{себе}.
\end{quote}

\begin{quote}
\emph{Mne ne khvatáyet uvérennosti v sebé.}
\end{quote}

\begin{quote}
I lack self-confidence.
\end{quote}

Feminine. \emph{\ru{Уверенность} \ru{в} \ru{себе}} — self-confidence. From \emph{\ru{уверенный}}, sure.

\begin{cousinweb}
\textbf{\ru{уверенность}} is built on \emph{\ru{уверенный}} \textquotedblleft{}sure\textquotedblright{}, from \emph{\ru{уверить}} \textquotedblleft{}to assure\textquotedblright{}, and so on \emph{\ru{вера}} \textquotedblleft{}faith\textquotedblright{}, inherited from Proto-Slavic \emph{*věra}; EtymDB names Latin \emph{vērus} \textquotedblleft{}true\textquotedblright{} as a cousin (EtymDB).
\end{cousinweb}

\begin{grammarlens}[title={What you've built}]
One more word for telling, explaining and describing at B1.
\end{grammarlens}

\subsection*{Your turn}
\begin{itemize}
\raggedright
  \item \emph{Say it:} \emph{uvérennost}
  \item \emph{Say it:} \emph{uvérennost}, once more
  \item \emph{Say it:} a sentence of your own with \emph{uvérennost}
  \item \emph{Recall:} say the Russian for to change someone’s mind, then the Russian for a conviction, then say \emph{uvérennost} again
\end{itemize}

\begin{tcolorbox}[breakable,colback=teal!4,colframe=teal!35!black,title={Before you move on}]
What does \ru{уверенность} mean? (\textquotedblleft{}Confidence, certainty\textquotedblright{}.) Say it once more.
\end{tcolorbox}
\section[\texorpdfstring{\ru{сомневаться} (somnevátsya)}{сомневаться (somnevátsya)}]{\ru{сомневаться} (somnevátsya) — to doubt}

\label{lesson:RU-C339-somnevatsya}



\begin{quote}
Before the new one: say the Russian for a conviction, then the Russian for confidence.
\end{quote}

\subsection*{You'll want to know: \ru{сомневаться}}
\textbf{\ru{сомневаться}} (\emph{somnevátsya}) — \textquotedblleft{}to doubt\textquotedblright{}.

\begin{quote}
\ru{Я} \ru{сомневаюсь}, \ru{что} \ru{он} \ru{придёт} \ru{вовремя}.
\end{quote}

\begin{quote}
\emph{Ya somneváyus, chto on pridyót vovrémya.}
\end{quote}

\begin{quote}
I doubt that he will come on time.
\end{quote}

Imperfective only. \emph{\ru{Сомневаться} \ru{в}} + prepositional: \emph{\ru{я} \ru{сомневаюсь} \ru{в} \ru{этом}}.

\begin{cousinweb}
\textbf{\ru{сомневаться}} is built on the root of \emph{\ru{мнить}} \textquotedblleft{}to think, to suppose\textquotedblright{} (EtymDB lists \emph{mnitĭ} \textquotedblleft{}to think\textquotedblright{}), the root of \emph{\ru{мнение}} \textquotedblleft{}opinion\textquotedblright{}, with the prefix \emph{\ru{со}-}. EtymDB sets Latin \emph{mēns} \textquotedblleft{}mind\textquotedblright{} beside \emph{\ru{мнить}} as a cousin.
\end{cousinweb}

\begin{grammarlens}[title={What you've built}]
One more word for telling, explaining and describing at B1.
\end{grammarlens}

\subsection*{Your turn}
\begin{itemize}
\raggedright
  \item \emph{Say it:} \emph{somnevátsya}
  \item \emph{Say it:} \emph{somnevátsya}, once more
  \item \emph{Say it:} a sentence of your own with \emph{somnevátsya}
  \item \emph{Recall:} say the Russian for a conviction, then the Russian for confidence, then say \emph{somnevátsya} again
\end{itemize}

\begin{tcolorbox}[breakable,colback=teal!4,colframe=teal!35!black,title={Before you move on}]
What does \ru{сомневаться} mean? (\textquotedblleft{}To doubt\textquotedblright{}.) Say it once more.
\end{tcolorbox}
